
\subsection{Modular Arithmetic \& Number Theory}

\begin{itemize}
    \item \textbf{Properties of Modulo:}
          \begin{itemize}
              \item Sum: $(a + b) \% m = (a \% m + b \% m) \% m$
              \item Subtraction: $(a - b) \% m = (a \% m - b \% m + m) \% m$
              \item Product: $(a \times b) \% m = (a \% m \times b \% m) \% m$
          \end{itemize}

    \item \textbf{Modular Division:}
          $(a / b) \equiv (a \times b^{-1}) \pmod m$, where $b \times b^{-1} \equiv 1 \pmod m$.
          Valid if $\gcd(b, m) = 1$.

    \item \textbf{Euler's Totient Theorem:}
          If $\gcd(a, m) = 1$, then $a^{\phi(m)} \equiv 1 \pmod m$.

    \item \textbf{Fermat's Little Theorem:}
          If $p$ is prime, $a^{p-1} \equiv 1 \pmod p$.
          Inverse: $a^{-1} \equiv a^{p-2} \pmod p$.

    \item \textbf{Power Reduction:}
          If $p$ is prime: $a^b \equiv a^{b \bmod (p-1)} \pmod p$.
          General: $a^b \equiv a^{b \bmod \phi(m)} \pmod m$ (if $\gcd(a, m) = 1$).

    \item \textbf{Bézout's Identity:}
          For $a, b$, there exist $x, y$ such that $ax + by = \gcd(a, b)$.
          If $\gcd(a, m) = 1$, use Extended Euclidean Algorithm to find $a^{-1} \pmod m$.

    \item \textbf{Chinese Remainder Theorem (CRT):}
          Given pairwise coprime moduli $m_1, \ldots, m_k$ and remainders $r_1, \ldots, r_k$, there exists a unique $x \pmod{M}$ (where $M = m_1 \cdots m_k$) such that $x \equiv r_i \pmod{m_i}$ for all $i$.
          \[
              x = \sum_{i=1}^{k} r_i \cdot \frac{M}{m_i} \cdot \left(\frac{M}{m_i}\right)^{-1}_{\!\!\!\!\mod m_i} \pmod M
          \]
          For two moduli: if $x \equiv a \pmod{m}$ and $x \equiv b \pmod{n}$ with $\gcd(m,n)=1$, then $x \equiv a + m \cdot ((b-a) \cdot m^{-1} \bmod n) \pmod{mn}$.
\end{itemize}
